%======================================================================
%  Capitulo 5 - Diseno fisico de bases de datos (generado por gen/cap5.py)
%======================================================================

\section{Diseño físico de bases de datos}

El diseño físico decide cómo se implementa el modelo relacional del
capítulo 4 en un sistema gestor de bases de datos concreto: tipos de
datos, generación de claves, índices, disparadores, usuarios y
permisos, dónde se guardan los datos y con qué recursos funciona la base
de datos. Como el resto del diseño detallado, se hace para el módulo
priorizado de operación de envíos (sección~\ref{sec:modulo}).

%-----------------------------------------------------------------------
\subsection{Conversión del Diseño Lógico al Diseño Físico}

\subsubsection{Modelo Relacional al Modelo Computacional (DBMS)}

\textbf{Sistema gestor de bases de datos.} Se usa Oracle Database,
uno de los gestores que indica el programa del curso (Oracle 23c o
superior), por tres razones: el modelo se construyó en Oracle SQL
Developer Data Modeler, que genera el esquema para Oracle; Oracle aplica todas las
restricciones del modelo lógico (claves, unicidad, CHECK y valores por
defecto) y permite programar en la base de datos las reglas que
dependen de otras filas; y tiene una edición gratuita, Oracle AI
Database Free, suficiente para la primera fase. Esa edición limita el
uso a dos núcleos de procesador, 2 GB de memoria y 12 GB de datos de
usuario\footnote{Oracle, \textit{Oracle AI Database Free Installation
Guide}, capítulo ``Licensing Restrictions'', 2026.}. La estimación de
la sección~\ref{sec:almacenamiento} muestra que el módulo ocupa cerca de
2,7 GB a los tres años, de modo que cabe holgadamente; cuando la
empresa supere esos límites se pasa a la edición Standard Edition 2,
que usa el mismo motor y no exige cambiar el diseño.

\textbf{Tipos de datos.} Los tipos del diccionario de datos
(tabla~\ref{tab:diccionario}) ya son tipos de Oracle; la
tabla~\ref{tab:tipos} resume el criterio con que se eligieron. La base de
datos usa el juego de caracteres AL32UTF8, y las longitudes de texto se
cuentan en caracteres y no en bytes, para que un nombre con tildes o
con eñe quepa en el tamaño definido.

{\singlespacing\small
\begin{longtable}{L{3.6cm} L{3.0cm} L{7.8cm}}
\caption{Correspondencia entre los datos del modelo y los tipos de Oracle.}
\label{tab:tipos}\\
\toprule
\textbf{Dato} & \textbf{Tipo en Oracle} & \textbf{Criterio} \\
\midrule
\endfirsthead
\toprule
\textbf{Dato} & \textbf{Tipo en Oracle} & \textbf{Criterio} \\
\midrule
\endhead
\bottomrule
\endfoot
Identificadores sustitutos & NUMBER(2) a NUMBER(12) & Tamaño según el volumen esperado: 2 dígitos para catálogos, 5 para centros, zonas y vehículos, 10 a 12 para las tablas que crecen a diario. \\
Número de guía & VARCHAR2(15) & Código alfanumérico impreso en la etiqueta: prefijo de la empresa y un consecutivo. \\
Nombres, direcciones y descripciones & VARCHAR2(n) & Longitud variable; solo ocupa el espacio del texto escrito. \\
Estados, tipos y listas de valores & VARCHAR2(10 a 25) & Valores legibles controlados con CHECK. \\
Indicadores sí/no & CHAR(1) & Valores S o N controlados con CHECK. \\
Pesos, capacidades y dinero & NUMBER(p,s) & Precisión fija: pesos con dos decimales en kg y valores con dos decimales en pesos colombianos. \\
Fechas de calendario & DATE & Fecha de registro, de compromiso, de vinculación. \\
Instantes de la operación & TIMESTAMP & Eventos, visitas y pruebas de entrega, donde importan la hora y el orden. \\
Soportes de la prueba de entrega & VARCHAR2(200) & Ruta del archivo (firma o fotografía); la imagen se guarda en el almacenamiento de archivos, no dentro de la base de datos. \\
\end{longtable}
\par}

\textbf{Generación de claves.} Las claves sustitutas se generan como
columnas de identidad, que Oracle implementa con una secuencia propia de
cada tabla; así ninguna aplicación tiene que calcular el siguiente
número. Las excepciones son tres. El número de guía lo asigna el
procedimiento de registro del envío, con el prefijo EE y un consecutivo
de 13 dígitos. La clave de MENSAJERO es la misma del empleado, porque es
un subtipo. Y los catálogos pequeños (ESTADO\_ENVIO, MODALIDAD y
TIPO\_VEHICULO) se cargan con valores fijos al instalar la base de
datos, porque las reglas y las vistas se refieren a ellos.

\textbf{Índices.} Oracle crea automáticamente un índice para cada clave
primaria y cada restricción de unicidad (32 índices), pero no para
las claves foráneas. Sin índice, cada consulta que une dos tablas por una
clave foránea, y cada verificación de que una fila no está referenciada,
recorre la tabla completa. Por eso se crea un índice para cada clave
foránea que no queda cubierta por un índice existente (22 índices)
y 5 índices adicionales para las consultas más frecuentes
(tabla~\ref{tab:indices}). En total el módulo tiene 59 índices.
Con el índice único sobre número de guía y fecha del evento, el estado
actual del envío se obtiene leyendo una sola entrada del índice, así que
no hace falta guardarlo como dato derivado en ENVIO (decisión 2 de la
optimización del diseño lógico).

{\singlespacing\small
\begin{longtable}{L{3.3cm} L{4.6cm} L{2.8cm} L{3.7cm}}
\caption{Índices que se agregan a los de las claves primarias y únicas.}
\label{tab:indices}\\
\toprule
\textbf{Tabla} & \textbf{Claves foráneas indexadas} & \textbf{Índice adicional} & \textbf{Motivo del índice adicional} \\
\midrule
\endfirsthead
\toprule
\textbf{Tabla} & \textbf{Claves foráneas indexadas} & \textbf{Índice adicional} & \textbf{Motivo del índice adicional} \\
\midrule
\endhead
\bottomrule
\endfoot
zona & id\_centro & -- &  \\
sede\_cliente & id\_cliente; id\_zona & -- &  \\
vehiculo & id\_centro; id\_tipo\_vehiculo & -- &  \\
empleado & id\_centro & -- &  \\
mensajero & id\_zona; id\_tipo\_vehiculo & -- &  \\
despacho & id\_centro; id\_vehiculo; id\_mensajero & fecha & Planilla del día y despachos abiertos. \\
solicitud\_recoleccion & id\_sede; id\_despacho & -- &  \\
envio & id\_solicitud; id\_destinatario; id\_zona; id\_modalidad & fecha\_registro & Indicadores y consultas por periodo. \\
asignacion\_despacho & id\_despacho & -- &  \\
evento\_envio & id\_estado; id\_empleado & numero\_guia, fecha\_hora (único) & Ubica el último evento de cada envío (estado actual y seguimiento) y evita dos eventos del mismo envío en el mismo instante; cubre también la clave foránea hacia ENVIO. \\
novedad & -- & estado & Novedades abiertas (vista de la mesa de servicio). \\
reclamacion & id\_cliente; numero\_guia & estado & Reclamaciones pendientes. \\
\end{longtable}
\par}

\textbf{Ciclo de vida del envío.} La regla RN12 exige que los estados de
un envío sigan un orden. La tabla~\ref{tab:transiciones} define las
transiciones permitidas; el disparador de EVENTO\_ENVIO rechaza
cualquier otra. ENTREGADO, DEVUELTO y ANULADO son estados finales: un
envío en uno de ellos no recibe más eventos.

{\singlespacing\small
\begin{longtable}{L{3.4cm} L{5.2cm} L{5.8cm}}
\caption{Transiciones permitidas entre estados del envío.}
\label{tab:transiciones}\\
\toprule
\textbf{Estado actual} & \textbf{Estados siguientes} & \textbf{Situación} \\
\midrule
\endfirsthead
\toprule
\textbf{Estado actual} & \textbf{Estados siguientes} & \textbf{Situación} \\
\midrule
\endhead
\bottomrule
\endfoot
(sin eventos) & REGISTRADO & Toda guía nace registrada (RN12). \\
REGISTRADO & RECOLECTADO, ANULADO & Se recoge en la sede o se anula antes de recogerlo. \\
RECOLECTADO & EN\_CENTRO, ANULADO & Llega al centro de distribución. \\
EN\_CENTRO & CLASIFICADO, ANULADO & Se pesa, se verifica y se clasifica por zona. \\
CLASIFICADO & EN\_RUTA, ANULADO & Sale en un despacho de entrega. \\
EN\_RUTA & ENTREGADO, CON\_NOVEDAD, DEVUELTO & Resultado del intento; DEVUELTO solo en el tercer intento (RN07). \\
CON\_NOVEDAD & EN\_RUTA & Reintento en un nuevo despacho. \\
ENTREGADO, DEVUELTO, ANULADO & -- & Estados finales. \\
\end{longtable}
\par}

\textbf{Disparadores y procedimientos.} Las reglas de negocio que
dependen de otras filas o tablas (tabla~\ref{tab:reglasnegocio}) se
programan en la base de datos, de modo que se cumplen sin importar qué
aplicación registre los datos. La tabla~\ref{tab:disparadores} muestra
qué hace cada uno; todos rechazan la operación con un mensaje que
explica la regla incumplida.

{\singlespacing\small
\begin{longtable}{L{4.0cm} L{3.4cm} L{5.4cm} L{1.4cm}}
\caption{Disparadores y procedimiento del módulo.}
\label{tab:disparadores}\\
\toprule
\textbf{Objeto} & \textbf{Tabla y momento} & \textbf{Qué verifica o hace} & \textbf{Reglas} \\
\midrule
\endfirsthead
\toprule
\textbf{Objeto} & \textbf{Tabla y momento} & \textbf{Qué verifica o hace} & \textbf{Reglas} \\
\midrule
\endhead
\bottomrule
\endfoot
Procedimiento de registro del envío & ENVIO y EVENTO\_ENVIO & Genera el número de guía, inserta el envío y su evento REGISTRADO en una sola transacción, y crea el destinatario si no existe. & RN02, RN12 \\
Disparador de solicitud & SOLICITUD\_RECOLECCION, antes de insertar & El cliente dueño de la sede está ACTIVO. & RN15 \\
Disparador de despacho & DESPACHO, antes de insertar o cambiar vehículo o mensajero & El vehículo está DISPONIBLE, el mensajero está ACTIVO y el vehículo es del tipo que conduce el mensajero. & RN10, RN11 \\
Disparador de asignación & ASIGNACION\_DESPACHO, antes de insertar & El despacho es de entrega y no está cerrado; el intento \textit{n} solo se crea si los anteriores terminaron en NOVEDAD; el tercero es la devolución. & RN07 \\
Disparador de resultado & ASIGNACION\_DESPACHO, antes de actualizar el resultado & ENTREGADO o DEVUELTO exigen una prueba de entrega del intento; NOVEDAD exige la novedad registrada; el tercer intento solo puede terminar en DEVUELTO. & RN05, RN07, RN08 \\
Disparador de prueba & PRUEBA\_ENTREGA, antes de insertar & Si el envío es frágil o de alto valor, el documento de quien recibe es el del destinatario. & RN06 \\
Disparador de evento & EVENTO\_ENVIO, antes de insertar & La transición está permitida (tabla~\ref{tab:transiciones}) y ENTREGADO o DEVUELTO corresponden a un intento con ese resultado. & RN05, RN12 \\
\end{longtable}
\par}

\textbf{Usuarios, roles y permisos.} Cada grupo de usuarios del módulo
tiene un rol con los permisos de la tabla~\ref{tab:permisos} traducidos a
privilegios: crear es insertar, actualizar es modificar y leer es
consultar. Ningún rol tiene permiso de borrar (RN16). Las tablas
pertenecen a un esquema propietario, BD\_ENLACE\_EXPRESS, con el que no
trabaja ningún usuario. El administrador de la base de datos crea los objetos,
carga los catálogos y, mientras se integran los módulos M2 a M4,
registra y actualiza los datos de referencia (clientes, sedes,
vehículos y empleados) a solicitud de las áreas dueñas
(tabla~\ref{tab:interfaces}). Los usuarios no tienen una cuenta propia
en la base de datos. El servidor de aplicaciones se conecta con una
cuenta de servicio por rol (cinco cuentas, como svc\_mesa y svc\_centro,
cada una con solo su rol) y elige la cuenta según el rol del usuario que
inició sesión en la aplicación; las consultas públicas del portal, que se
hacen antes de iniciar sesión, usan la cuenta de menor privilegio, la del
rol cliente. En cada conexión, el servidor fija además el identificador
de cliente de la sesión de Oracle con el usuario real, de modo que la
base de datos sabe quién hizo cada operación y puede auditarla. Así no
hay que crear ni administrar en la base de datos una cuenta por cada
empleado y cliente: el control por usuario lo hace la aplicación (por
ejemplo, el mensajero solo ve sus despachos y el cliente solo sus
envíos), y Oracle sigue limitando, mediante la cuenta de servicio, lo que
cada rol puede hacer.

{\singlespacing\small
\begin{longtable}{L{2.3cm} L{0.9cm} L{3.6cm} L{3.6cm} L{3.5cm}}
\caption{Roles de la base de datos y sus privilegios.}
\label{tab:roles}\\
\toprule
\textbf{Rol} & \textbf{Usu.} & \textbf{Inserta} & \textbf{Modifica} & \textbf{Consulta} \\
\midrule
\endfirsthead
\toprule
\textbf{Rol} & \textbf{Usu.} & \textbf{Inserta} & \textbf{Modifica} & \textbf{Consulta} \\
\midrule
\endhead
\bottomrule
\endfoot
rol\_gerencia & U1 & -- & -- & Tablas que alimentan los indicadores y vistas del módulo \\
rol\_centro & U2 & centro\_distribucion, zona, despacho, asignacion\_despacho, evento\_envio & centro\_distribucion, zona, despacho, asignacion\_despacho, solicitud\_recoleccion, envio, empleado, mensajero & Catálogos, sedes, vehículos, empleados; vistas de planilla y envíos por despachar \\
rol\_mesa & U3 & solicitud\_recoleccion, envio, destinatario, evento\_envio, reclamacion & envio, destinatario, novedad, reclamacion & Clientes, sedes, zonas, catálogos; vistas de seguimiento, novedades y reclamaciones \\
rol\_mensajero & U4 & evento\_envio, prueba\_entrega, novedad & solicitud\_recoleccion, asignacion\_despacho, envio & Sedes de recolección, su despacho y la vista de planilla \\
rol\_cliente & U9 & -- & -- & Vista de seguimiento, filtrada por el portal \\
\end{longtable}
\par}

\textbf{Espacios de tablas.} Las tablas y los índices se guardan en
espacios de tablas separados (tabla~\ref{tab:tablespaces}), lo que
permite vigilar el crecimiento de cada uno y, si se cambia de servidor,
ubicarlos en discos distintos. Su tamaño inicial cubre el primer año de
la estimación y crecen automáticamente hasta un máximo que respeta el
límite de 12 GB de la edición gratuita.

{\singlespacing\small
\begin{longtable}{L{3.6cm} L{3.8cm} L{1.6cm} L{2.2cm} L{1.6cm}}
\caption{Espacios de tablas del módulo.}
\label{tab:tablespaces}\\
\toprule
\textbf{Espacio de tablas} & \textbf{Contenido} & \textbf{Inicial} & \textbf{Aumento} & \textbf{Máximo} \\
\midrule
\endfirsthead
\toprule
\textbf{Espacio de tablas} & \textbf{Contenido} & \textbf{Inicial} & \textbf{Aumento} & \textbf{Máximo} \\
\midrule
\endhead
\bottomrule
\endfoot
TS\_ENLACE\_DATOS & Las 19 tablas & 500 MB & 100 MB & 7 GB \\
TS\_ENLACE\_INDICES & Los 59 índices & 300 MB & 100 MB & 4 GB \\
\end{longtable}
\par}

%-----------------------------------------------------------------------
\subsubsection{Organización física de los datos}

\paragraph{Almacenamiento (Centralizado o Distribuido)}
\label{sec:almacenamiento}

El almacenamiento es \textbf{centralizado}: una sola base de datos
guarda la información de todos los centros. Se descarta el
almacenamiento distribuido porque la empresa opera en una sola ciudad;
un envío pasa por varios centros y zonas y su historial debe estar
completo en un solo lugar; la política de sistema único desde el inicio
(capítulo 1) busca precisamente no repartir los datos; y el volumen
esperado lo maneja sin dificultad un solo servidor. Distribuir los datos
obligaría a replicarlos y a mantener la consistencia entre sitios sin
ningún beneficio para este tamaño de operación.

Para dimensionar el almacenamiento se partió de los datos del
capítulo 1 con los siguientes supuestos: 900 envíos por día hábil
(25 mensajeros con unas 36 visitas cada uno), 300 días hábiles al año
(lunes a sábado) y un crecimiento del 25\,\% anual. Por cada envío se
registran en promedio 6,2 eventos, 1,12 asignaciones a despacho, una
prueba de entrega y 0,12 novedades; el 60\,\% de los envíos va a un
destinatario nuevo, una solicitud de recolección agrupa unos seis envíos
y cada mensajero sale en dos despachos al día. El tamaño de cada fila se
calculó con los tipos del diccionario de datos, suponiendo los textos
llenos a la mitad, y el de los bloques con un 10\,\% de espacio libre
para actualizaciones. La tabla~\ref{tab:volumen} muestra el resultado
para las tablas que crecen con la operación; las demás son tablas
maestras y de catálogo de pocos cientos de filas.

{\singlespacing\small
\begin{longtable}{L{3.4cm} >{\raggedleft\arraybackslash}p{1.4cm} >{\raggedleft\arraybackslash}p{2.0cm} >{\raggedleft\arraybackslash}p{2.0cm} >{\raggedleft\arraybackslash}p{2.0cm} >{\raggedleft\arraybackslash}p{2.0cm}}
\caption{Estimación del volumen de las tablas que crecen con la operación.}
\label{tab:volumen}\\
\toprule
\textbf{Tabla} & \textbf{Bytes por fila} & \textbf{Filas año 1} & \textbf{Filas a 3 años} & \textbf{Datos a 3 años (MB)} & \textbf{Índices a 3 años (MB)} \\
\midrule
\endfirsthead
\toprule
\textbf{Tabla} & \textbf{Bytes por fila} & \textbf{Filas año 1} & \textbf{Filas a 3 años} & \textbf{Datos a 3 años (MB)} & \textbf{Índices a 3 años (MB)} \\
\midrule
\endhead
\bottomrule
\endfoot
evento\_envio & 144 & 1.674.000 & 6.382.125 & 987,2 & 685,6 \\
envio & 234 & 270.000 & 1.029.375 & 258,7 & 137,4 \\
prueba\_entrega & 198 & 270.000 & 1.029.375 & 218,9 & 49,0 \\
destinatario & 136 & 162.000 & 617.625 & 90,2 & 32,2 \\
asignacion\_despacho & 54 & 302.400 & 1.152.900 & 66,9 & 129,2 \\
novedad & 348 & 32.400 & 123.525 & 46,2 & 8,5 \\
solicitud\_recoleccion & 55 & 45.000 & 171.562 & 10,1 & 11,8 \\
reclamacion & 567 & 2.700 & 10.294 & 6,3 & 0,9 \\
despacho & 74 & 15.000 & 57.188 & 4,5 & 6,3 \\
\midrule
\textbf{Total (19 tablas)} & & & & \textbf{1689,3} & \textbf{1061,0} \\
\end{longtable}
\par}

La base de datos del módulo ocupa unos 0,7 GB al final del primer
año y 2,7 GB a los tres años, incluidos los índices. EVENTO\_ENVIO es
la tabla más grande, porque cada envío deja varios eventos. Las
fotografías y firmas de las pruebas de entrega (unos 200 KB cada una,
cerca de 51 GB al año) se guardan en un almacenamiento de archivos
aparte, y la base de datos solo conserva su ruta.

\textbf{Copias de seguridad.} La base de datos trabaja en modo
ARCHIVELOG, que conserva el registro de todos los cambios. Se hace una
copia completa semanal y una incremental diaria con RMAN, la
herramienta de copias de Oracle, y se copian los registros de cambios
cada hora; las copias se guardan fuera del servidor durante 30 días.
Con ello se puede recuperar la base de datos hasta la última hora de
operación. Además, cada mes se exporta el esquema con Data Pump como
respaldo lógico.

\paragraph{Procesamiento (Centralizado o Distribuido)}

El procesamiento también es \textbf{centralizado}, en una arquitectura
cliente-servidor de tres niveles físicos. Los usuarios trabajan en
aplicaciones cliente (web para el centro, la mesa de servicio y la
gerencia; móvil para los mensajeros; portal para los clientes); estas
se comunican con un servidor de aplicaciones que expone servicios web,
y solo ese servidor se conecta a la base de datos. Dentro del servidor
de aplicaciones, el software se organiza en capas lógicas (casos de
uso, dominio y acceso a datos), como se justifica en la
sección~\ref{sec:arquitectura}. Las reglas de negocio se verifican en
la capa de dominio y, como garantía final, en la base de datos
(restricciones, disparadores y procedimientos), así que se cumplen
igual sin importar qué aplicación registre la operación.

La carga es baja para un solo servidor. Con los supuestos anteriores,
la operación genera unas 9.000 filas nuevas por día, la mayoría
eventos; si el 20\,\% se concentra en la hora de mayor actividad, son
menos de una inserción por segundo, con unos 35 usuarios internos
conectados al tiempo más las consultas del portal. Como los mensajeros
pueden quedarse sin señal en la ruta, la aplicación móvil guarda los
registros en el teléfono y los envía al recuperar la conexión, con la
hora real del evento.

%-----------------------------------------------------------------------
\subsubsection{Recursos Hardware, Software, Aplicaciones y BD, Comunicaciones, Ingeniería Web}

\textbf{Hardware.} La tabla~\ref{tab:hardware} resume los equipos. Los
servidores se contratan como máquinas virtuales en la nube, lo que evita
comprar equipos para una empresa que apenas inicia y permite ampliarlos
cuando crezca la operación.

{\singlespacing\small
\begin{longtable}{L{3.4cm} L{1.4cm} L{9.6cm}}
\caption{Recursos de hardware.}
\label{tab:hardware}\\
\toprule
\textbf{Recurso} & \textbf{Cant.} & \textbf{Características y uso} \\
\midrule
\endfirsthead
\toprule
\textbf{Recurso} & \textbf{Cant.} & \textbf{Características y uso} \\
\midrule
\endhead
\bottomrule
\endfoot
Servidor de base de datos & 1 & Máquina virtual con 2 núcleos, 8 GB de memoria y 100 GB de disco de estado sólido (base de datos, registros de cambios y copias locales). \\
Servidor de aplicaciones & 1 & Máquina virtual con 2 núcleos y 4 GB de memoria; atiende las aplicaciones web, la móvil y el portal. \\
Almacenamiento de archivos & 1 & Almacenamiento en la nube para fotografías y firmas de las pruebas de entrega y para las copias de seguridad (unos 51 GB al año más las copias). \\
Computadores de escritorio & 10 & Puestos del personal del centro, la mesa de servicio y la gerencia. \\
Teléfonos inteligentes & 25 & Uno por mensajero, con cámara, GPS y plan de datos, para la aplicación móvil. \\
Lectores de código de barras & 3 & Lectura de la guía al recibir, clasificar y despachar los paquetes en el centro. \\
Impresoras de etiquetas & 2 & Impresión de la etiqueta con el número de guía. \\
Equipos de red & 1 & Enrutador con cortafuegos en el centro de distribución. \\
\end{longtable}
\par}

\textbf{Software.} La tabla~\ref{tab:software} lista el software de la
base de datos y del servidor.

{\singlespacing\small
\begin{longtable}{L{3.6cm} L{4.6cm} L{6.2cm}}
\caption{Recursos de software.}
\label{tab:software}\\
\toprule
\textbf{Función} & \textbf{Software} & \textbf{Uso} \\
\midrule
\endfirsthead
\toprule
\textbf{Función} & \textbf{Software} & \textbf{Uso} \\
\midrule
\endhead
\bottomrule
\endfoot
Gestor de bases de datos & Oracle AI Database Free (después, Standard Edition 2) & Motor de la base de datos del módulo. \\
Sistema operativo & Oracle Linux & Sistema de los dos servidores. \\
Modelado & Oracle SQL Developer Data Modeler & Modelos conceptual, lógico y físico; diccionario de datos. \\
Administración y consultas & Oracle SQL Developer & Creación de objetos, consultas, informes y administración. \\
Servidor de aplicaciones & Aplicación propia con acceso a datos por JDBC & Implementa las capas de casos de uso, dominio y datos, y expone los servicios web que usan las aplicaciones (sección~\ref{sec:arquitectura}). \\
Copias de seguridad & RMAN y Data Pump & Copias físicas y exportaciones lógicas; vienen con la base de datos. \\
\end{longtable}
\par}

\textbf{Aplicaciones y base de datos.} Cada aplicación atiende a un grupo
de usuarios y se conecta con el rol de ese grupo
(tabla~\ref{tab:aplicaciones}). Las aplicaciones completas no forman parte del alcance del proyecto
(sección~\ref{sec:modulo}); el prototipo funcional
(sección~\ref{sec:prototipo}) demuestra sus funciones principales sobre
la base de datos del módulo en una sola aplicación web, con una sección
para cada grupo de usuarios y el portal público de rastreo.

{\singlespacing\small
\begin{longtable}{L{3.6cm} L{1.0cm} L{2.6cm} L{7.2cm}}
\caption{Aplicaciones que usan la base de datos.}
\label{tab:aplicaciones}\\
\toprule
\textbf{Aplicación} & \textbf{Usu.} & \textbf{Rol} & \textbf{Funciones principales} \\
\midrule
\endfirsthead
\toprule
\textbf{Aplicación} & \textbf{Usu.} & \textbf{Rol} & \textbf{Funciones principales} \\
\midrule
\endhead
\bottomrule
\endfoot
Aplicación de operación (web) & U2 & rol\_centro & Centros y zonas, recepción y clasificación, armado y cierre de despachos. \\
Aplicación de mesa de servicio (web) & U3 & rol\_mesa & Solicitudes, registro de envíos, novedades, consultas y reclamaciones. \\
Aplicación del mensajero (móvil) & U4 & rol\_mensajero & Planilla del día, recolecciones, eventos, pruebas de entrega con foto o firma y novedades; funciona sin conexión. \\
Tablero de indicadores (web) & U1 & rol\_gerencia & Indicadores de cumplimiento, novedades y reclamaciones. \\
Portal de seguimiento (web) & U9 & rol\_cliente & Consulta del estado y el recorrido de los envíos del cliente. \\
\end{longtable}
\par}

\textbf{Comunicaciones.} El centro de distribución tiene una conexión a
internet de fibra óptica y una conexión móvil de respaldo; los
mensajeros usan los datos móviles del teléfono. Todo el tráfico entre
las aplicaciones y el servidor de aplicaciones viaja cifrado (HTTPS). La
base de datos no se expone a internet: está en una red privada a la que
solo llegan el servidor de aplicaciones y, para administración, los
equipos autorizados mediante una red privada virtual (VPN).

\textbf{Ingeniería web.} Las aplicaciones web y el portal son
adaptables a la pantalla (computador o teléfono) y consumen los
servicios web del servidor de aplicaciones en formato JSON. Cada usuario
inicia sesión con su cuenta, y la aplicación solo le muestra las
opciones de su rol. El portal de clientes es la única parte pública:
permite consultar un envío con su número de guía y, con inicio de
sesión, todos los envíos del cliente. Los servicios registran quién
hizo cada operación, lo que complementa el responsable que ya guarda
cada evento del envío.
